\documentclass[12pt]{article}
\usepackage{amsmath, amssymb}
\usepackage{geometry}
\usepackage{enumitem}
\usepackage{tcolorbox}
\usepackage{xcolor}
\usepackage{booktabs}

\geometry{margin=1in}

\definecolor{answerblue}{RGB}{230, 240, 255}
\definecolor{questiongray}{RGB}{245, 245, 245}

\tcbset{
  answerbox/.style={
    colback=answerblue,
    colframe=blue!40,
    fonttitle=\bfseries,
    title=Answer,
    rounded corners,
    boxrule=0.5pt
  },
  questionbox/.style={
    colback=questiongray,
    colframe=gray!40,
    rounded corners,
    boxrule=0.5pt
  }
}

\title{\textbf{ESE 650: Learning in Robotics}\\
\large HMM \& Viterbi --- Practice Problems}
\author{}
\date{}

\begin{document}
\maketitle

\section*{Part 1: HMM Core Concepts (Q1--Q10)}

%% Q1
\subsection*{Q1}
\begin{tcolorbox}[questionbox]
\textbf{True or False:} The forward variable $\alpha_k(x) = P(X_k = x \mid Y_1, \ldots, Y_k)$,
i.e., it is the posterior distribution over states given observations.
\end{tcolorbox}

\begin{tcolorbox}[answerbox]
\textbf{False.} The forward variable is a \emph{joint} probability, not a conditional:
\[
  \alpha_k(x) = P(X_k = x,\, Y_1, \ldots, Y_k).
\]
To obtain the filtering distribution, normalize over all states:
\[
  P(X_k = x \mid Y_{1:k}) = \frac{\alpha_k(x)}{\displaystyle\sum_{x'} \alpha_k(x')}.
\]
\end{tcolorbox}

%% Q2
\subsection*{Q2}
\begin{tcolorbox}[questionbox]
\textbf{True or False:} The Viterbi algorithm and the forward algorithm have the same
time complexity, $O(T \cdot |X|^2)$.
\end{tcolorbox}

\begin{tcolorbox}[answerbox]
\textbf{True.} Both algorithms share the same recursive structure. At each of the $T$
time steps, every pair $(x', x)$ of consecutive states must be considered, giving
$|X|^2$ operations per step and $O(T \cdot |X|^2)$ overall. The only difference is
that the forward algorithm uses \textbf{sum} while Viterbi uses \textbf{max}.
\end{tcolorbox}

%% Q3
\subsection*{Q3}
\begin{tcolorbox}[questionbox]
\textbf{True or False:} The smoothing estimate $P(X_k \mid Y_1, \ldots, Y_T)$ always
has lower entropy than the filtering estimate $P(X_k \mid Y_1, \ldots, Y_k)$.
\end{tcolorbox}

\begin{tcolorbox}[answerbox]
\textbf{False.} Smoothing uses additional future observations and is therefore
\emph{generally} more certain, but not \emph{always}. If the future observations
$Y_{k+1}, \ldots, Y_T$ carry no information about $X_k$ (e.g.\ the emission matrix
row for every state is uniform), the two distributions are identical and their
entropies are equal. The correct statement is:
\[
  H\!\left(P(X_k \mid Y_{1:T})\right) \;\leq\; H\!\left(P(X_k \mid Y_{1:k})\right).
\]
\end{tcolorbox}

%% Q4
\subsection*{Q4}
\begin{tcolorbox}[questionbox]
Standard Viterbi returns only the most likely path. What modification is needed
to also return the \emph{probability} of that path?
\end{tcolorbox}

\begin{tcolorbox}[answerbox]
\textbf{No modification is needed.} The Viterbi variable $\delta_k(x)$ already stores
the probability of the most likely path ending in state $x$ at time $k$. After
back-tracking to recover the optimal path $x_1^*, \ldots, x_T^*$, the probability
is simply $\delta_T(x_T^*)$.
\end{tcolorbox}

%% Q5
\subsection*{Q5}
\begin{tcolorbox}[questionbox]
Suppose row $x^*$ of the emission matrix $M$ is uniform, i.e.\
$M_{x^*, y} = \text{const}$ for all $y$. What is the effect on (a) filtering
and (b) Viterbi decoding?
\end{tcolorbox}

\begin{tcolorbox}[answerbox]
In both cases the observation \emph{cannot distinguish} state $x^*$ from other
states.

\begin{itemize}
  \item \textbf{Filtering:} The likelihood term for state $x^*$ is the same
        constant for every observation, so incoming observations never update
        the relative probability of $x^*$. Its estimate evolves purely through
        the transition probabilities.
  \item \textbf{Viterbi:} The emission factor for any path passing through $x^*$
        is constant across all time steps and observations. Whether $x^*$ appears
        in the optimal path depends entirely on the transition probabilities, not
        on the observations.
\end{itemize}
\end{tcolorbox}

%% Q6
\subsection*{Q6}
\begin{tcolorbox}[questionbox]
Describe how to compute the joint smoothing probability of two consecutive states:
\[
  P(X_k = a,\, X_{k+1} = b \mid Y_1, \ldots, Y_T).
\]
\end{tcolorbox}

\begin{tcolorbox}[answerbox]
Using the forward--backward decomposition:
\[
  P(X_k = a, X_{k+1} = b \mid Y_{1:T})
  \;\propto\;
  \underbrace{\alpha_k(a)}_{\text{past}}
  \cdot
  \underbrace{T_{ab}}_{\text{transition}}
  \cdot
  \underbrace{M_{b,\, y_{k+1}}}_{\text{emission}}
  \cdot
  \underbrace{\beta_{k+1}(b)}_{\text{future}},
\]
where $\beta_k(x) = P(Y_{k+1}, \ldots, Y_T \mid X_k = x)$ is the backward variable.
Normalize by dividing by $P(Y_{1:T}) = \sum_{x} \alpha_T(x)$.
\end{tcolorbox}

%% Q7
\subsection*{Q7}
\begin{tcolorbox}[questionbox]
A student claims: ``I run the forward algorithm to get $\alpha_k(x)$, then normalize
it to obtain the filtering result $P(X_k \mid Y_1, \ldots, Y_k)$.'' Is this correct?
\end{tcolorbox}

\begin{tcolorbox}[answerbox]
\textbf{True.} Since $\alpha_k(x) = P(X_k = x, Y_{1:k})$, normalizing over all
states yields exactly the Bayes filter output:
\[
  P(X_k = x \mid Y_{1:k}) = \frac{\alpha_k(x)}{\sum_{x'} \alpha_k(x')}.
\]
\end{tcolorbox}

%% Q8
\subsection*{Q8}
\begin{tcolorbox}[questionbox]
\textbf{True or False:} If some observation $Y_t$ has very small emission probability
under all states, the forward variables will suffer from numerical underflow.
How is this addressed in practice?
\end{tcolorbox}

\begin{tcolorbox}[answerbox]
\textbf{True.} The forward recursion multiplies emission probabilities at every
step. Repeated multiplication of small values causes floating-point underflow to
zero. The standard fix is to work in \textbf{log-space}, converting products into
sums:
\[
  \log \alpha_k(x) = \log \alpha_{k-1}(x') + \log T_{x'x} + \log M_{x,\, y_k}.
\]
\end{tcolorbox}

%% Q9
\subsection*{Q9}
\begin{tcolorbox}[questionbox]
Consider an HMM with a very large state space ($|X| = 10^6$) such that storing
the full trellis is infeasible. Can you still perform (a) filtering and
(b) smoothing? Explain.
\end{tcolorbox}

\begin{tcolorbox}[answerbox]
\begin{itemize}
  \item \textbf{Filtering: Yes.} The forward recursion only requires the previous
        time step's $\alpha_{k-1}(\cdot)$ to compute $\alpha_k(\cdot)$. Memory
        usage is $O(|X|)$; past steps can be discarded.
  \item \textbf{Smoothing: Very difficult.} The forward--backward algorithm
        requires $\alpha_k(x)$ and $\beta_k(x)$ at the \emph{same} time step.
        Since $\beta$ is computed backwards, all $T$ forward vectors must be
        stored simultaneously, requiring $O(T \cdot |X|)$ memory --- infeasible
        for $|X| = 10^6$.
\end{itemize}
\end{tcolorbox}

%% Q10
\subsection*{Q10}
\begin{tcolorbox}[questionbox]
We sample hallucinated future observations $\hat{Y}_{k+1}, \ldots, \hat{Y}_{k+T}$
and use them to compute a smoothing estimate of $X_k$. If we repeat this process
infinitely many times and average the results, what does the average converge to?
\end{tcolorbox}

\begin{tcolorbox}[answerbox]
The average converges to the \textbf{standard filtering estimate}:
\[
  P(X_k \mid Y_{1:k})
  = \mathbb{E}_{\hat{Y}_{k+1},\ldots,\hat{Y}_{k+T}}
    \!\Big[P\!\left(X_k \mid Y_{1:k},\,\hat{Y}_{k+1},\ldots,\hat{Y}_{k+T}\right)\Big].
\]
Marginalizing over all possible hallucinated futures cancels their effect,
regardless of how they were sampled. This also confirms (Q3) that a single
hallucinated future does \emph{not} improve the estimate.
\end{tcolorbox}

%% ============================================================
\newpage
\section*{Part 2: Viterbi Deep Dive (VQ1--VQ3)}

%% VQ1
\subsection*{VQ1}
\begin{tcolorbox}[questionbox]
What algorithm do you obtain if you replace every $\max$ in Viterbi with $\sum$?
What is the key difference in what each algorithm computes?
\end{tcolorbox}

\begin{tcolorbox}[answerbox]
Replacing $\max$ with $\sum$ yields the \textbf{forward algorithm}. The two
recursions are:
\[
  \text{Forward:}\quad
  \alpha_{k+1}(x) = \sum_{x'} \alpha_k(x')\, T_{x'x}\, M_{x,\, y_{k+1}},
\]
\[
  \text{Viterbi:}\quad
  \delta_{k+1}(x) = \max_{x'} \delta_k(x')\, T_{x'x}\, M_{x,\, y_{k+1}}.
\]
The forward algorithm \textbf{sums over all paths} to compute the marginal
likelihood $P(Y_{1:k}, X_k = x)$, while Viterbi \textbf{retains only the best
path} probability up to each state.
\end{tcolorbox}

%% VQ2
\subsection*{VQ2}
\begin{tcolorbox}[questionbox]
Viterbi finds the path maximizing the \emph{joint} probability
$P(X_{1:T}, Y_{1:T})$. An alternative approach picks the most likely state
at each time step independently:
\[
  \hat{x}_k = \argmax_{x}\; P(X_k = x \mid Y_{1:T}).
\]
Do these two methods always agree?
\end{tcolorbox}

\begin{tcolorbox}[answerbox]
\textbf{No.} The per-step argmax ignores transition probabilities between time
steps, so the resulting sequence $(\hat{x}_1, \ldots, \hat{x}_T)$ may be
\emph{globally inconsistent} --- e.g., it may contain transitions
$\hat{x}_k \to \hat{x}_{k+1}$ with near-zero probability under the model.

Viterbi enforces \textbf{global consistency} by jointly optimizing over the
entire path. The per-step approach (obtained via forward--backward smoothing)
is useful when per-marginal accuracy matters more than path coherence.
\end{tcolorbox}

%% VQ3
\subsection*{VQ3}
\begin{tcolorbox}[questionbox]
Can Viterbi be applied directly to a \emph{continuous} state space
(e.g., robot position $x \in \mathbb{R}^2$)? If not, what are the alternatives?
\end{tcolorbox}

\begin{tcolorbox}[answerbox]
\textbf{No.} Viterbi requires computing $\max_{x' \in \mathcal{X}}$ at each step.
For a continuous state space this maximization is over infinitely many values and
cannot be enumerated.

Two practical alternatives:
\begin{itemize}
  \item \textbf{Particle Viterbi:} Approximate the continuous space with a finite
        set of particles. Once discretized, the standard Viterbi $\max$ can be
        applied over the particle set.
  \item \textbf{Kalman smoother (RTS smoother):} For linear--Gaussian systems,
        the optimal state sequence can be computed analytically. The forward pass
        is the Kalman filter; the backward pass is the Rauch--Tung--Striebel
        smoother.
\end{itemize}
\end{tcolorbox}

\end{document}
